\documentclass[10pt]{beamer}
\usepackage[utf8]{inputenc}
\usepackage[spanish]{babel}
\usepackage{booktabs}
\usepackage{hyperref}

% ==========================================================
% PERSONALIZACIÓN DE LA PLANTILLA
% ==========================================================
\usetheme{Madrid}

% 1. Personalización de Paleta de Colores (Colores Institucionales PUCP)
\definecolor{pucpNavy}{RGB}{0, 43, 73}
\definecolor{pucpGold}{RGB}{197, 155, 39}

\setbeamercolor{structure}{fg=pucpNavy}
\setbeamercolor{palette primary}{bg=pucpNavy, fg=white}
\setbeamercolor{palette secondary}{bg=pucpNavy!85, fg=white}
\setbeamercolor{palette tertiary}{bg=pucpNavy!70, fg=white}
\setbeamercolor{titlelike}{parent=palette primary}
\setbeamercolor{block title}{bg=pucpNavy, fg=white}
\setbeamercolor{block body}{bg=pucpNavy!10, fg=black}
\setbeamercolor{block title alert}{bg=pucpGold!90!black, fg=white}
\setbeamercolor{block body alert}{bg=pucpGold!15, fg=black}

% 2. Personalización Completa del Pie de Página (Footer Structure)
\setbeamertemplate{footline}{
  \leavevmode%
  \hbox{%
  \begin{beamercolorbox}[wd=.333333\paperwidth,ht=2.25ex,dp=1ex,center]{author in head/foot}%
    \usebeamerfont{author in head/foot}\insertshortauthor
  \end{beamercolorbox}%
  \begin{beamercolorbox}[wd=.333333\paperwidth,ht=2.25ex,dp=1ex,center]{title in head/foot}%
    \usebeamerfont{title in head/foot}Economía PUCP
  \end{beamercolorbox}%
  \begin{beamercolorbox}[wd=.333333\paperwidth,ht=2.25ex,dp=1ex,right]{date in head/foot}%
    \usebeamerfont{date in head/foot}\insertshortdate{}\hspace*{2em}
    \insertframenumber{} / \inserttotalframenumber\hspace*{2ex}
  \end{beamercolorbox}}%
  \vskip0pt%
}
\setbeamertemplate{navigation symbols}{}

% --- DATOS DEL DOCUMENTO ---
\title[Crecimiento e Inflación]{Análisis del Crecimiento Económico y la Inflación en el Perú}
\subtitle{Evaluación Macroeconómica Reciente}
\author[Estudiante PUCP]{Estudiante de Economía}
\institute[PUCP]{Pontificia Universidad Católica del Perú}
\date{\today}

\begin{document}

% SLIDE 1: Carátula
\begin{frame}
  \titlepage
\end{frame}

% SLIDE 2: Agenda
\begin{frame}{Estructura de la Presentación}
  \tableofcontents
\end{frame}

\section{Introducción y Marco Teórico}

% SLIDE 3: Bloques de Énfasis
\begin{frame}{Motivación del Estudio}
  \begin{block}{Objetivo Principal}
    Analizar los determinantes del crecimiento económico peruano y su interacción con las metas de inflación del BCRP.
  \end{block}

  \vspace{0.4cm}

  \begin{alertblock}{Diagnóstico Central}
    La estabilidad macroeconómica es condición indispensable pero insuficiente para garantizar el desarrollo inclusivo a largo plazo.
  \end{alertblock}
\end{frame}

% SLIDE 4: Revelación Progresiva (Ritmo de Composición)
\begin{frame}{Modelo Neoclásico de Solow}
  Supuestos fundamentales del modelo de crecimiento:
  \begin{itemize}[<+->]
    \item Rendimientos decrecientes en la acumulación del capital.
    \item Función de producción agregada: $Y = A K^\alpha L^{1-\alpha}$.
    \item Tasa de ahorro constante $s$ y depreciación del capital $\delta$.
    \item Progreso tecnológico exógeno $g$.
  \end{itemize}
\end{frame}

\section{Análisis de Datos y Resultados}

% SLIDE 5: Entorno Columns (Composición Avanzada)
\begin{frame}{Desempeño Macroeconómico Reciente}
  \begin{columns}[t]
    \begin{column}{0.48\textwidth}
      \begin{block}{Crecimiento del PBI}
        \begin{itemize}
          \item Rebote histórico post-pandemia en 2021 (+13.4\%).
          \item Desaceleración estructural en 2022--2023.
          \item Urgencia de mayor productividad.
        \end{itemize}
      \end{block}
    \end{column}

    \begin{column}{0.48\textwidth}
      \begin{alertblock}{Control Inflacionario}
        \begin{itemize}
          \item Choques globales de oferta en 2022.
          \item Incremento preventivo de la tasa de referencia.
          \item Convergencia al rango meta (1\%--3\%).
        \end{itemize}
      \end{alertblock}
    \end{column}
  \end{columns}
\end{frame}

% SLIDE 6: Tabla
\begin{frame}{Datos Macroeconómicos Consolidados}
  \centering
  \small
  \begin{tabular}{lccc}
    \toprule
    \textbf{Indicador} & \textbf{2021} & \textbf{2022} & \textbf{2023} \\
    \midrule
    Crecimiento PBI (\%) & 13.4 & 2.7 & -0.6 \\
    Inflación Anual (\%) & 6.4 & 8.5 & 3.2 \\
    Tasa Referencia BCRP (\%) & 2.50 & 7.50 & 6.75 \\
    \bottomrule
  \end{tabular}
\end{frame}

% SLIDE 7: Revelación Progresiva y Botón Interactivo (Beamer Button)
\begin{frame}[label=slide_monetaria]{Mecanismos de Política Monetaria}
  Pasos del canal de transmisión de la tasa de interés:

  \begin{enumerate}[<+->]
    \item Ajuste de la tasa de interés de referencia por el BCRP.
    \item Transmisión al costo del crédito del sistema financiero.
    \item Modulación de la demanda agregada e inflación.
  \end{enumerate}

  \vspace{0.6cm}
  
  \centering
  % Botón hipervinculado al Apéndice
  \hyperlink{slide_apendice}{\beamerbutton{Ver Derivación Matemática en Apéndice}}
\end{frame}

\section{Políticas Económicas y Conclusiones}

% SLIDE 8: Columnas para Análisis
\begin{frame}{Retos y Oportunidades Estructurales}
  \begin{columns}
    \begin{column}{0.5\textwidth}
      \begin{block}{Retos Estructurales}
        \begin{itemize}
          \item Alta informalidad laboral (>70\%).
          \item Brecha de infraestructura pública.
        \end{itemize}
      \end{block}
    \end{column}
    \begin{column}{0.5\textwidth}
      \begin{block}{Oportunidades}
        \begin{itemize}
          \item Sólida posición de reservas (RIN).
          \item Alta demanda de minerales de transición.
        \end{itemize}
      \end{block}
    \end{column}
  \end{columns}
\end{frame}

% SLIDE 9: Conclusiones
\begin{frame}{Conclusiones}
  \begin{itemize}[<+->]
    \item La estabilidad de precios preserva el poder adquisitivo del sol.
    \item Se requiere impulsar la inversión privada para retomar tasas de crecimiento sostenidas superiores al 3\%.
    \item La institucionalidad del BCRP es la fortaleza principal del modelo macroeconómico peruano.
  \end{itemize}
\end{frame}

% SLIDE 10: Cierre
\begin{frame}{Cierre}
  \centering
  \Large \textbf{¡Muchas Gracias!}

  \vspace{1cm}
  \normalsize ¿Preguntas o Comentarios?
\end{frame}

% SLIDE 11: Apéndice (Destino Interactivo del Botón)
\begin{frame}[label=slide_apendice]{Apéndice: Ecuaciones del Estado Estacionario}
  \begin{block}{Ecuación de Acumulación}
    $$\dot{k} = s f(k) - (n + g + \delta)k$$
  \end{block}

  En estado estacionario ($\dot{k} = 0$):
  $$s k^*{}^\alpha = (n + g + \delta)k^* \implies k^* = \left( \frac{s}{n + g + \delta} \right)^{\frac{1}{1-\alpha}}$$

  \vspace{0.5cm}
  \centering
  \hyperlink{slide_monetaria}{\beamerbutton{Volver a la Presentación}}
\end{frame}

\end{document}
